%%
%% spring-lecture04.tex
%% 
%% Made by Alex Nelson <pqnelson@gmail.com>
%% Login   <alex@lisp>
%% 
%% Started on  2026-04-09T11:35:38-0700
%% Last update 2026-04-09T11:35:38-0700
%% 

\lecture{}

\begin{definition}
A \define{Complete Lattice} consists of a partially ordered set
(``poset'') such that
\begin{enumerate}
\item every subset has a greatest lower bound, and
\item every subset has a least upper bound.
\end{enumerate}
\end{definition}

\begin{node}
For an $A$-module $M$, we write $S(M)$ for the set of all submodules
of $M$. Then $S(M)$ is a complete lattice. Thie poset structure is
given by inclusion. For any $X=\{N_{i}\}_{i\in J}\subset S(M)$, we see
$\bigcap_{j\in J}N_{j}$ is the greatest lower bound and $\sum_{j\in J}N_{j}$
is the least upper bound.

Furthermore, $S(M)$ satisfies the \define{Modular Law} (or ``modularity''):
for any $N,P,Q\in S(M)$, if $N\subset Q$, then $N+(P\cap Q)=(N+P)\cap Q$.
\end{node}

\begin{proof}
We see $N+(P\cap Q)\subset N+P$ and $N+(P\cap Q)\subset N+Q\subset Q$,
hence $N+(P\cap Q)\subset(N+P)\cap Q$.

For any $u=v+w\in(N+P)\cap Q$ where $v\in N$ and $w\in P$ we see $u\in Q$,
then $w=u-v\in P\cap Q$, and $u=v+w\in N+(P\cap Q)$.
Hence $N+(P\cap Q)\supset(N+P)\cap Q$.
\end{proof}

\begin{lemma}\label{lemma:2.7}
For an $A$-module $M$, the following are equivalent:
\begin{enumerate}
\item For any $N,P,Q\in S(M)$, $N\cap(P+Q)=(N\cap P)+(N\cap Q)$
\item For any $N,P,Q\in S(M)$, $N+(P\cap Q)=(N+P)\cap(N+Q)$
\item For any $N,P,Q\in S(M)$, $(N\cap P)+(P\cap Q)+(Q\cap N)=(N+P)\cap(P+Q)\cap(Q+N)$.
\end{enumerate}
If $S(M)$ satisfies any of these equivalent conditions, we say $S(M)$
is \define{Distributive}.
\end{lemma}

\begin{proof}
$(1)\implies(2)$ By using (1) on the underlined items (the first line
  matches the left-hand side of (1) to give us the second line, then the second line matches the
  right-hand side of (1) to give us the third line):
\begin{subequations}
  \begin{align}
\underline{(N+P)}\cap(\underline{N}+\underline{Q}) &= \bigl((N+P)\cap N\bigr)+\bigl((N+P)\cap Q\bigr)\\
&=\underline{((N\cap N)+(P\cap N))}+\underline{((N\cap Q)+(P\cap Q))}\\
&=N+(P\cap Q).
  \end{align}
\end{subequations}
Hence (2).

$(2)\implies(3)$ We apply (2) to the underlined items:
\begin{subequations}
  \begin{align}
\underline{(N\cap P)}+(\underline{P}\cap\underline{Q})+(Q\cap N)
&=\bigl(((N\cap P)+P)\cap((N\cap P)+Q)\bigr)+(Q\cap N)\\
&=\bigl((N+P)\cap(P+P)\bigr)\cap((N+Q)\cap(P+Q))\;+\;(Q\cap N)\\
&=P\cap(N+Q)\cap(P+Q)\;+\;(Q\cap N)\\\intertext{but since $P\cap(N+Q)\supset(P+Q)$,}
&=P\cap(N+Q)\;+\;(Q\cap N)\\
&=(P+(Q\cap N))\cap(N+Q)\\
&=(P+Q)\cap(P+N)\cap(N+Q)
  \end{align}
\end{subequations}
by applying (2) twice. Hence (3).

$(3)\implies(1)$
First we apply modularity to get
\begin{subequations}
  \begin{align}
N\cap(P+Q) &=N\cap\underline{(N+P)\cap(P+Q)\cap(Q+N)}\\
&=N\cap\bigl(\underline{(N\cap P)}+\underline{(P\cap Q)+(Q\cap N)}\bigr)\mbox{ by (3)}\\
&=(N\cap P)+(Q\cap N)+(N\cap P\cap Q)\mbox{ by modularity}\\
&=(N\cap P)+(Q\cap N),
  \end{align}
\end{subequations}
since $N\cap P\cap Q\subset Q\cap N$ and $N\cap P\cap Q\subset N\cap P$.
Hence (1).
\end{proof}

\begin{proposition}\label{prop:2.8}
Let $M$ be an $A$-module such that $S(M)$ is not distributive. Then
there exists some $P\neq Q\in S(M)$ such that
\begin{equation}
\frac{P}{P\cap Q}\iso\frac{Q}{P\cap Q}
\end{equation}
as $A$-modules.
\end{proposition}

\begin{proof}
Since $S(M)$ is not distributive, there exists $M_{1},M_{2},M_{3}\in S(M)$
such that
\begin{subequations}
\begin{equation}
M_{0}:=(M_{1}\cap M_{2})+(M_{2}\cap M_{3})+(M_{3}\cap M_{1})
\end{equation}
\begin{equation}
M_{4}:=(M_{1}+M_{2})\cap(M_{2}+M_{3})\cap(M_{3}+M_{1})
\end{equation}
and
\begin{equation}
M_{0}\propersubset M_{4}.
\end{equation}
\end{subequations}
Each summand in $M_{0}$ is properly contained in each, uh,
``intersectand'' of $M_{4}$. Let
\begin{subequations}
  \begin{align}
N &:= (M_{1}\cap M_{2})+(M_{3}\cap(M_{1}+M_{2}))\\
P &:= (M_{2}\cap M_{3})+(M_{1}\cap(M_{2}+M_{3}))\\
Q &:= (M_{3}\cap M_{1})+(M_{2}\cap(M_{3}+M_{1}))
  \end{align}
\end{subequations}
(this is $(M_{i}\cap M_{j})+(M_{k}\cap(M_{i}+M_{j}))$ where $(i,j,k)$
is $(1,2,3)$, $(2,3,1)$, $(3,1,2)$ respectively).
Then by repeatedly applying the modular laws to underlined terms, we have
\begin{subequations}
  \begin{align}
P\cap Q 
&= \underline{[(M_{2}\cap M_{3})+(M_{1}\cap(M_{2}+M_{3}))]}\cap[\underline{(M_{3}\cap M_{1})}+\underline{(M_{2}\cap(M_{3}+M_{1}))}]\\
&= (M_{3}\cap M_{1})+\biggl(\underline{(M_{2}\cap M_{3})}+\bigl(\underline{(M_{1}\cap(M_{2}+M_{3}))}\cap\underline{(M_{2}\cap(M_{3}+M_{1}))}\bigl)\biggl)\\
&=(M_{3}\cap M_{1})+(M_{2}\cap M_{3})+\bigl(M_{1}\cap(M_{2}+M_{3})\cap M_{2}\cap(M_{3}+M_{1})\bigr)\\
&=(M_{3}\cap M_{1})+(M_{2}\cap M_{3})+(M_{1}\cap M_{2})\\
&=M_{0}.
  \end{align}
\end{subequations}
We also see since $M_{2}\cap M_{3}\subset M_{2}\cap(M_{3}+M_{1})$ and
$M_{3}\cap M_{1}\subset M_{1}\cap(M_{2}+M_{3})$
\begin{subequations}
  \begin{align}
P+Q&=(M_{2}\cap M_{3})+(M_{1}\cap(M_{2}+M_{3}))+(M_{3}\cap M_{1})+(M_{2}\cap(M_{3}+M_{1}))\\
&=(M_{1}\cap(M_{2}+M_{3}))+(M_{2}\cap(M_{3}+M_{1}))\\
&=\bigl((M_{1}\cap(M_{2}+M_{3}))+M_{2}\bigr)\cap(M_{1}+M_{3})\\
&=m_{4}.
  \end{align}
\end{subequations}
Similarly,
\begin{equation}
P\cap N=Q\cap N=M_{0}
\end{equation}
and
\begin{equation}
P+N=Q+N=M_{4}.
\end{equation}
In particular, $P\neq Q$, and if we take
\begin{subequations}
  \begin{align}
\frac{P}{P\cap Q} &\iso\frac{P}{P\cap N}\\
&\iso\frac{(P+N)}{N}\\
&\iso\frac{M_{4}}{N}.
  \end{align}
\end{subequations}
The same chain of isomorphisms gives us
\begin{subequations}
  \begin{align}
\frac{Q}{P\cap Q}&\iso\frac{Q}{Q\cap N}\\
&\iso\frac{Q+N}{N}\\
&\iso\frac{M_{4}}{N}.
  \end{align}
\end{subequations}
Hence the claim.
\end{proof}

\begin{definition}
A right or left $A$-module is called \define{Simple} (or
\define{Irreducible}) if $M\neq0$ and $S(M)=\{0,M\}$.

If $M$ is the direct sum o simple $A$-modules, then we say $M$ is
\define{Semisimple} (resp., \define{Completely Irreducible}).
\end{definition}

\begin{proposition}\label{prop:2.10}
For an $A$-module $N$, the following are equivalent:
\begin{enumerate}
\item $N$ is simple
\item For any nonzero element $u\in N$, we have $N=uA$
\item $N\iso A/M$ for some maximal right ideal $M$ of $A$
\end{enumerate}
\end{proposition}

\begin{proof}
$(1)\implies(2)$ since $uA\neq0$, then $uA=N$.

$(2)\implies(1)$ If $uA=N$ for every $u\neq0$, then $N$ has no
nontrivialsubmodules.

$(3)\implies(1)$ Follows from the correspondence theorem (i.e., there
exists a bijection between $S(A/M)$ and $\{M'\in S(A)\mid M\subset M'\}$).

$(2)\implies(3)$ Let $0\neq u\in N$. The map $A\to N$, $x\mapsto ux$
is surjective and its kernel $M$ is a right ideal. Since $N\iso A/M$i
s simply (by (1)), by the correspondence theorem $M$ is maximal.
\end{proof}

\begin{lemma}[Schur's lemma]\label{lemma:2.11}
Let $\phi\colon M\to N$ be a nonzero $A$-module morphism. If $M$
(resp., $N$) is simple, then $\phi$ is injective (resp., surjective).

In particular, if $M$ and $N$ are simple, then either $M\iso N$ or
there is no such $\phi$ (i.e., $\hom_{A}(M,N)=\emptyset$).
\end{lemma}

\begin{proof}
Since $\phi\neq0$, $\ker(\phi)\neq M$ and $\Im(\phi)\neq0$. By the
simplicity of $M\cap N$ the claim is proved.
\end{proof}

\begin{corollary}\label{cor:2.12}
For a \emph{semisimple} $A$-module $N$, the following are equivalent:
\begin{enumerate}
\item $N$ is simple
\item $\End_{A}(N)$ is a division algebra
\item $N$ is \define{Indecomposable} [i.e., $N\neq0$ and there are no
  $P,Q\in S(M)\setminus\{0\}$ such that $N=P\oplus Q$].
\end{enumerate}
\end{corollary}

\begin{proof}
$(1)\implies(2)$ By Schur's lemma, any $\phi\in\End_{A}(N)$ is
  bijective hence invertible.

$(2)\implies(3)$ If $N=P\oplus Q$, then there exists some $\pi\in\End_{A}(N)$
such that $\pi(N)=P$ [so $(\id-\pi)(N)=Q$] and $\pi^{2}=\pi$. Hence
$\pi(\id-\pi)=0$. But then either $\pi=0$ or $\id-\pi=0$, which
implies either $P=0$ or $Q=0$ respectively.

$(3)\implies(1)$ Trivial.
\end{proof}